\documentclass[12pt]{article}

\usepackage{sbc-template}
\usepackage{graphicx,url}
\usepackage[utf8]{inputenc}
\usepackage[T1]{fontenc}
\usepackage[brazil]{babel}
\usepackage{amsmath,amssymb}
\usepackage{booktabs}
\usepackage[hidelinks]{hyperref}

\sloppy

\title{Mineração de Métricas DORA em Repositórios Open-Source\\ que Utilizam GitHub Actions}

\author{Guilherme Martini Brina Ferreira\inst{1}\\
        Gabriel Cardoso\inst{1}}

\address{Pontifícia Universidade Católica de Minas Gerais (PUC Minas)\\
  Belo Horizonte -- MG -- Brasil
  \email{gmbferreira12@gmail.com, gabriel2004cardoso@gmail.com}
}

\begin{document}

\maketitle

\begin{abstract}
The DevOps Research and Assessment (DORA) metrics have become the de facto
standard for measuring software delivery performance. This work investigates
whether those metrics, originally designed for organizations, can be inferred
from the public activity of open-source projects that adopt GitHub Actions. We
mine releases, workflow runs and commits over a twelve-month window to compute
deployment frequency, lead time for changes, change failure rate and failed
deployment recovery time, plus a sensitivity analysis of the operational
definitions. This paper presents the context, the seven research questions and
the informal hypotheses that guide the study before the data are analysed.
\end{abstract}

\begin{resumo}
As métricas do \emph{DevOps Research and Assessment} (DORA) tornaram-se o padrão
de fato para medir o desempenho de entrega de software. Este trabalho investiga
se essas métricas, originalmente pensadas para organizações, podem ser
inferidas a partir da atividade pública de projetos open-source que adotam o
GitHub Actions. Mineramos releases, execuções de workflow e commits ao longo de
uma janela de doze meses para calcular frequência de deploy, lead time for
changes, taxa de falha de mudança e tempo de recuperação, além de uma análise de
sensibilidade das definições operacionais. Este artigo apresenta o contexto, as
sete questões de pesquisa e as hipóteses informais que orientam o estudo antes
da análise dos dados.
\end{resumo}

\section{Introdução}

As métricas DORA (\emph{DevOps Research and Assessment}) consolidaram-se como a
principal referência de mercado para medir o desempenho de entrega de software.
Popularizadas pelo livro \emph{Accelerate} \cite{forsgren2018} e atualizadas
anualmente pelo relatório \emph{Accelerate State of DevOps} \cite{dorastate}, elas
resumem o desempenho de uma equipe em quatro dimensões: duas de
\emph{velocidade} --- a frequência de deploy e o tempo de espera das mudanças
(\emph{lead time for changes}) --- e duas de \emph{estabilidade} --- a taxa de
falha de mudança (\emph{change failure rate}) e o tempo de recuperação após uma
falha \cite{doraguide}. A definição dessas métricas é revisada com o tempo: o
antigo \emph{mean time to restore} (MTTR) passou a ser chamado de
\emph{failed deployment recovery time} em 2023 e, em 2024, o DORA acrescentou
uma quinta métrica, a \emph{deployment rework rate} \cite{dorahistory}.

Essas métricas foram originalmente elaboradas a partir de questionários
respondidos por profissionais de empresas, e não de dados públicos de
repositórios. O ecossistema open-source, no entanto, tornou-se um objeto de
estudo atrativo por sua escala e por sua transparência: ferramentas como o
GitHub Actions expõem, de forma aberta e auditável, a cadência de integração
contínua, o histórico de releases e as mensagens de commit de milhares de
projetos. Surge então uma oportunidade relevante: é possível \emph{minerar}
métricas DORA a partir dessa atividade pública, em vez de depender de
auto-relato?

A resposta exige cautela. O GitHub \textbf{não registra} ``deploys em produção''
nem ``falhas em produção''. Um projeto open-source publica \emph{releases}, roda
\emph{workflows} de CI e recebe \emph{commits}, e é apenas a partir desses
eventos que as métricas DORA precisam ser \emph{inferidas}. Uma medida indireta
como essa é um \emph{proxy}: ``release publicada'' representa ``deploy'' e
``execução de CI que falhou'' representa ``falha''. Proxies podem enganar. Uma
biblioteca que publica uma release por mês pode não estar ``deployando'' nada em
produção --- quem implanta são seus usuários. Medir um proxy sem questioná-lo
ameaça a validade de construto do estudo \cite{wohlin2012}.

Neste trabalho, portanto, a pergunta não é apenas ``qual é o valor da métrica?'',
mas também ``o quanto podemos confiar nesse valor?''. Para responder a isso,
mineramos, via API REST do GitHub, releases publicadas, execuções de workflow
disparadas por \texttt{push} no \emph{default branch} e os commits contidos entre
releases, em uma janela de observação de doze meses. A partir desses dados
calculamos as métricas DORA em variantes operacionais obrigatórias e comparamos
como a classificação de desempenho de um repositório muda quando a definição
escolhida muda.

O estudo é guiado por sete questões de pesquisa. As quatro primeiras medem,
respectivamente, a frequência de deploy \textbf{(RQ~01)}, o tempo entre commit e
deploy \textbf{(RQ~02)}, a taxa de falha das mudanças entregues
\textbf{(RQ~03)} e o tempo de recuperação após uma falha \textbf{(RQ~04)}. A
\textbf{RQ~05} investiga a relação entre velocidade e estabilidade, a
\textbf{RQ~06} busca características de repositório associadas a melhor
desempenho e a \textbf{RQ~07} quantifica o quanto as conclusões dependem da
definição operacional adotada (\emph{análise de sensibilidade}). As hipóteses
informais que orientam cada questão, formuladas antes de qualquer análise dos
dados, são apresentadas na Seção~\ref{sec:hipoteses}.

As contribuições esperadas são: (i) um pipeline próprio de coleta e cálculo das
métricas DORA para projetos open-source; (ii) um retrato descritivo do
desempenho de entrega de repositórios populares sob definições operacionais
explícitas e comparáveis; e (iii) uma avaliação da robustez dessas conclusões
frente a escolhas metodológicas alternativas. O restante do artigo está
organizado da seguinte forma: a Seção~\ref{sec:hipoteses} apresenta as questões
de pesquisa e as hipóteses; a Metodologia, os Resultados e a Discussão são
desenvolvidos nas próximas etapas do projeto.

\section{Questões de Pesquisa e Hipóteses}
\label{sec:hipoteses}

Esta seção enuncia cada questão de pesquisa acompanhada da hipótese informal do
grupo --- o que se espera encontrar e por quê --- registrada \emph{antes} de
observar os resultados. As hipóteses são qualitativas e não comprometem a
análise: servem de referência para a discussão posterior (hipótese
\emph{versus} resultado) e para explicitar possíveis vieses dos autores.

\subsection*{RQ 01 --- Qual a frequência de deploys dos repositórios populares que usam CI/CD?}

\textbf{Hipótese.} Esperamos uma distribuição fortemente assimétrica, com a
mediana na faixa \emph{High} (ao menos um deploy por semana e menos de sete),
tipicamente entre uma e três releases por semana. Uma minoria de projetos
mantém cadência quase diária (\emph{Elite}, sete ou mais por semana), graças à
automação de release, enquanto um grupo pequeno publica menos de uma release por
mês (\emph{Low}). A justificativa é que projetos populares têm processos de
entrega maduros, mas a maioria não automatiza totalmente a publicação, e muitos
seguem um ciclo de releases mais conservador.

\subsection*{RQ 02 --- Qual o tempo entre um commit e seu respectivo deploy?}

\textbf{Hipótese.} A variante por release (a) produzirá valores bem maiores que
a variante por commit (b): esperamos a mediana de (a) na casa de semanas a um
mês e a de (b) na casa de dias. A razão é estrutural: na variante (a), um único
commit antigo ``esquecido'' em uma release faz o lead time daquela release
explodir, e esse valor extremo empurra a mediana para cima em projetos com
histórico irregular; na variante (b), cada commit contribui individualmente, de
modo que poucos valores extremos diluem-se no conjunto. Consequentemente,
muitos repositórios devem migrar para categorias melhores sob a variante (b),
e a divergência entre (a) e (b) é, ela própria, um resultado a ser discutido.

\subsection*{RQ 03 --- Qual a taxa de falha das mudanças entregues por esses repositórios?}

\textbf{Hipótese.} As duas variantes devem divergir bastante. O proxy de CI (a)
--- a proporção de execuções de workflow que falharam --- deve ser
comparativamente alto, com mediana em torno de 20\% a 40\% (faixa
\emph{Medium}), porque pipelines falham por motivos variados (testes, linters,
\emph{flakiness}) que não chegam a afetar o usuário. Já o proxy de entrega (b)
--- releases seguidas de correção em até sete dias --- deve ser baixo, com
mediana igual ou inferior a 15\% (\emph{Elite}/\emph{High}), pois releases
corretivas rápidas são relativamente raras e nem toda correção sai como nova
release. Esperamos que a diferença entre (a) e (b) seja um dos achados centrais
sobre a distinção entre falha de pipeline e falha de entrega.

\subsection*{RQ 04 --- Qual o tempo de recuperação após uma execução de CI/CD com falha?}

\textbf{Hipótese.} Esperamos medianas curtas, inferiores a um dia e
provavelmente inferiores a uma hora (faixa \emph{Elite} ou \emph{High}). Uma
execução de CI que falha costuma ser corrigida por uma re-execução ou um
\emph{hotfix} de pipeline no mesmo dia, já que o custo de corrigir um pipeline
quebrado é baixo. Ainda assim, prevemos uma cauda relevante de episódios
\emph{censurados} --- falhas que não foram resolvidas antes do fim da janela ---
sobretudo em projetos com manutenção intermitente; a proporção de censura deve
ser reportada por repositório.

\subsection*{RQ 05 --- Repositórios com maior frequência de deploy apresentam maior ou menor taxa de falha?}

\textbf{Hipótese.} Esperamos uma correlação de Spearman fraca e próxima de zero;
se houver significância, inclinamo-nos a uma associação levemente negativa (mais
deploys associados a menor CFR). Essa expectativa é alinhada à tese do DORA de
que velocidade e estabilidade \emph{não} formam um \emph{trade-off}
\cite{forsgren2018}. Não descartamos, contudo, que a variante de CFR afete o
resultado: a correlação com o proxy de CI (a) pode ser diferente da obtida com o
proxy de entrega (b).

\subsection*{RQ 06 --- Quais características dos repositórios estão associadas a um melhor desempenho DORA?}

\textbf{Hipótese.} Entre os fatores analisados (linguagem, popularidade,
número de contribuidores, idade e tipo de projeto), acreditamos que o
\textbf{tipo de projeto} e a \textbf{idade} sejam os mais associados ao
desempenho. Bibliotecas e \emph{frameworks} tendem a ter releases mais
espaçadas e \emph{lead time} maior, enquanto aplicações e serviços --- com
entrega contínua --- tendem a maior frequência de deploy; repositórios mais
antigos devem apresentar pior desempenho, por acúmulo de dívida técnica e base
de código maior. Esperamos que a \textbf{linguagem principal} tenha efeito
pequeno e que, após a correção de Holm para comparações múltiplas, a maioria
dos efeitos significativos seja de tamanho pequeno a médio.

\subsection*{RQ 07 --- O quanto a classificação DORA de um repositório depende da definição operacional escolhida?}

\textbf{Hipótese.} A classificação deve mostrar-se sensível sobretudo à
\textbf{unidade de deploy}: alternar entre ``release'', ``release +
pré-release'' e ``tag'' tende a mover uma parcela relevante (estimamos 20\% a
40\%) dos repositórios de categoria. A escolha do \emph{lead time} entre as
variantes (a) e (b) também desloca repositórios, enquanto a troca do proxy de
CFR deve ter efeito menor na categoria \emph{geral}. Esperamos concordância
global \emph{moderada} (kappa ponderado entre 0,4 e 0,6), indicando que as
conclusões das RQ~01 a RQ~06 são apenas parcialmente robustas e que a definição
operacional precisa ser explicitada em qualquer comparação.

% ---------------------------------------------------------------------------
% Link do repositório e do GitHub Projects do grupo (exigido pela disciplina).
% ---------------------------------------------------------------------------
\section*{Link do Repositório e do GitHub Projects}

\noindent Repositório: \url{https://github.com/Gmbferreira/Mineracao-Metricas-DORA}\\
GitHub Projects: \url{https://github.com/users/Gmbferreira/projects/2/views/1}

% ---------------------------------------------------------------------------
% Seções a serem desenvolvidas nas próximas sprints (Lab03S02 e Lab03S03).
% Mantidas aqui como esqueleto, propositalmente vazias nesta entrega.
% ---------------------------------------------------------------------------
% \section{Metodologia}
%   \subsection{Fonte de dados e funil de seleção}
%   \subsection{Janela de observação e definições operacionais}
%   \subsection{Variantes obrigatórias}
%   \subsection{Protocolo de validação manual}
%
% \section{Resultados}
%   \subsection{RQ 01 a RQ 04 e classificação DORA de referência}
%   \subsection{RQ 05 e RQ 06}
%   \subsection{RQ 07 --- análise de sensibilidade}
%
% \section{Discussão}
%
% \section{Ameaças à Validade}
%
% \section{Replicação Cruzada}

\bibliographystyle{sbc}
\bibliography{referencias}

\end{document}
